% Author: Grayson Orr
% Course: ID607001: Introductory Application Development Concept

\documentclass{article}
\author{}

\usepackage{graphicx}
\usepackage{wrapfig}
\usepackage{enumerate}
\usepackage{hyperref}
\usepackage[margin = 2.25cm]{geometry}
\usepackage[table]{xcolor}
\usepackage{fancyhdr}
\hypersetup{
  colorlinks = true, 
  urlcolor = blue
}
\setlength\parindent{0pt}
\pagestyle{fancy}
\fancyhf{}
\rhead{College of Engineering, Construction \& Living Sciences\\Bachelor of Information Technology}
\lfoot{In-Class Activity: ES6 Basics Extension\\Version 1, Semester Two, 2021}
\rfoot{\thepage}

\begin{document}

\begin{figure}
  \centering
  \includegraphics[width=50mm]{../img/logo.png}
\end{figure}

\title{College of Engineering, Construction \& Living Sciences\\Bachelor of Information Technology\\ID607001: Introductory Application Development Concepts\\Level 6, Credits 15\\\textbf{In-Class Activity: ES6 Basics Extension}}
\date{}
\maketitle

\section*{Instructions}
The purpose of this in-class activity is to extend yourself with more complex problems. Also, you will articulate various \textbf{JavaScript} concepts.

\section*{Submission}
You must submit all program files via \textbf{GitHub Classroom}. Here is the URL to the repository you will use for your code review – \href{https://classroom.github.com/a/\_6KSahyX}{https://classroom.github.com/a/\_6KSahyX}. If you wish to have your code reviewed, message the course lecturer on \textbf{Microsoft Teams}.

\section*{Getting Started}
Open your repository in \textbf{Visual Studio Code}. Create a new file called \textbf{03-in-class-activity.js}. In \textbf{03-in-class-activity.js}, add the following: 

\begin{verbatim}
  console.log('Hello, World!')
\end{verbatim}

Open a \textbf{terminal} \& run the following command:

\begin{verbatim}
  node 03-in-class-activity.js
\end{verbatim}

If the output is \textbf{Hello, World!}, then you are ready to start coding.

\subsection*{Problem 1:}
You have been given two \textbf{arrays} containing the lecturer's favourite programming languages. Use the following hints to display the expected output:
\begin{itemize}
  \item Add a specified element to the end of a list.
  \item Add all elements of a specified \textbf{array} to the end of a list.
  \item If present, remove a specified element from a \textbf{array}.
  \item Capitalise the element in the 3rd index.
\end{itemize}

\begin{verbatim}
  const progLangsOne = ['C#', 'JavaScript', 'Kotlin', 'OCaml']
  const progLangsTwo = ['C++', 'Go', 'Swift', 'TypeScript']

  // Write your solution here

  // Expected output:
  // [C#, JavaScript, Kotlin, OCAML, Prolog, C++, Swift]
\end{verbatim}

\subsection*{Problem 2:}
Create an \textbf{arrow function} which simulates the \textbf{Rock, Paper, Scissors} game. The \textbf{arrow function} takes the input of two players (rock, paper or scissors), first parameter from the first player, second from the second player. The function returns the result as such:

\begin{itemize}
  \item First player wins
  \item Second player wins
  \item Draw
\end{itemize}

\begin{verbatim}
  // Write your solution here

  console.log(rockPaperScissor('paper', 'rock'))
  console.log(rockPaperScissor('rock', 'paper'))
  console.log(rockPaperScissor('paper', 'paper'))

  // Expected output:
  // First player wins
  // Second player wins
  // Draw
\end{verbatim}

\subsection*{Problem 3:}
You are given the length of a song in minutes. The format is \textbf{mm:ss}, i.e., minutes:seconds or "01:00". Create an \textbf{arrow function} that takes the song's length \& returns it in seconds. \textbf{Note:} If the number of seconds is ≥ 60, return \textbf{false}.

\begin{verbatim}
  // Write your solution here

  console.log(minToSecs('01:00'))
  console.log(minToSecs('13:56'))
  console.log(minToSecs('10:60'))

  // Expected output:
  // 60
  // 836
  // false
\end{verbatim}

\subsection*{Problem 4:}
Explain the difference between \textbf{const} \& \textbf{Object.freeze()} using the following example:

\begin{verbatim}
  // const example
  const fruit = { name: 'Apple' }
  let vege = { name: 'Carrot' }
  fruit = vege

  // Object.freeze() example
  let person = { name: 'John' }
  Object.freeze(person)
  person.name = 'Jane'
  console.log(person)
\end{verbatim}

\subsection*{Problem 5:}
Explain the difference between \textbf{implicit} \& \textbf{explicit} type coercion.

\subsection*{Problem 6:} 
Explain what an \textbf{Immediately Invoked Function Expression (IIFE)} is. Provide an example in your answer.

\subsection*{Problem 7:}
Explain what a \textbf{generator} is. Provide an example in your answer.

\subsection*{Problem 8:}
Explain difference between an \textbf{undeclared}, \textbf{undefined} \& \textbf{null} variable.

\subsection*{Problem 9:}
Provide \textbf{three} examples on how you can empty an \textbf{array}.

\subsection*{Problem 10:}
Explain the expected output for the following example:

\begin{verbatim}
  const progLangsOne = ['C#', 'JavaScript', 'Kotlin', 'OCaml']
  delete progLangsOne[2]
  console.log(progLangsOne)
\end{verbatim}

\subsection*{Problem 11:}
Explain the operators used \& expected output in the following example:

\begin{verbatim}
  +true
  !'John'
\end{verbatim}

\subsection*{Problem 12:}
Explain the difference between \textbf{prefix} \& \textbf{postfix} increment/decrement operators.

\subsection*{Problem 13:}
Explain what \textbf{currying} is using the following example:

\begin{verbatim}
  const makeSandwich = (ingredientOne) => {
      return (ingredientTwo) => {
          return (ingredientThree) => {
              return `${ingredientOne}, ${ingredientTwo}, ${ingredientThree}`
          }
      }
  }

  const sandwich = makeSandwich('Bacon')('Lettuce')('Tomato')
  console.log(sandwich) // Bacon, Lettuce, Tomato
\end{verbatim}

\subsection*{Problem 14:}
Explain the following \textbf{expressions} \& \textbf{operators}:
\begin{itemize}
  \item Exponentiation assignment (**=)
  \item Logical nullish assignment (??=)
  \item Nullish coalescing operator (??)
  \item Optional chaining (?.)
  \item Remainder assignment (\%=)
\end{itemize}

\subsection*{Problem 15:}
Explain what a \textbf{Set} is. Provide an example in your answer.

\end{document}
